\silentchapter{Tre brev}{tre brev}
\begin{widebody}
\begin{foreordbody}

\vspace{4mm}
\noindent{\scshape\addfontfeature{LetterSpace=6}\small til designaren}
\vspace{2mm}

Du har lese hit, og du veit allereie det meste. Du kjenner eingongskoppen, overvakingsarkitekturen, merksemdsøkonomien. Du har sett Palantir. Du har sett kva faget tillèt deg å kalle godt. Diagnosen i etterord II er ikkje ny for deg; ho er berre sjeldan namngjeven.

Det denne boka tilbyr er ikkje ein ny metode. Det er eit vokabular. Med det kan du seie kva du gjer, kva du utelèt, og kvifor; ikkje som rettferdiggjering, men som rekneskap. Ingen formgjevar arbeider utan aksar som vert valde bort. Skiljet mellom det faget i dag har og det faget kan ha, ligg i om denne utelatinga vert skjult eller synleg. Synleggjering er starten på ein profesjon.

Du treng ikkje slutte i oppdrag for å byrje. Du kan byrje i neste leveranse. Dokumenter omhugen. Namngje kven som ber kostnaden. Det er nok i fyrste steg.

\vspace{8mm}
\noindent{\scshape\addfontfeature{LetterSpace=6}\small til studenten}
\vspace{2mm}

Du står ved inngangen. Pensumet du får er under bygging; det du vel å krevje i løpet av studiet, vil forme kva faget vert. Krev at kvart prosjekt inneheld ei skriftleg grunngjeving for aksane som vart representerte og dei som vart valde bort. Krev at handverket er substansielt; at du får tid i materialet, ikkje berre i representasjonen. Krev at nabofaga vert konsulterte når forma du arbeider med rører ved levande system, ved store mengder menneske, ved økologi, ved makt.

Vel oppdrag etter kva dei krev av deg, ikkje etter kva dei betalar. Ein kandidat som aldri har sagt nei til eit oppdrag har ikkje ferdigstilt utdanninga. Ho har ferdigstilt ei kopiering. Å seie nei er eit profesjonelt grep, ikkje eit moralsk. Utan det er profesjonen ikkje ein profesjon.

Du har ein fordel dei som er midt i karrieren ikkje har: du kan byrje rett. Det er òg eit ansvar.

\vspace{8mm}
\noindent{\scshape\addfontfeature{LetterSpace=6}\small til akademiet}
\vspace{2mm}

Rammeverket er falsifiserbart; vilkåra er lista i notata. Fundament-seksjonen viser at same formalismen er demonstrert over fjorten substratnivå. Forskingsprogrammet er fire punkt: operasjonell definisjon av kjerneomgrepa, kvantifisering av omhug gjennom K-metrikken, reanalysering av eksisterande arkiv, sertifisering bunden til dokumentert breidde.

Det som manglar er ikkje teori; det er institusjonell forpliktelse. Eit organ som kan ekskludere for brot, slik legeforeininga kan. Eit pensum som bind handverket til substrat og nabofag. Ei publiserinngsinfrastruktur som aksepterer at formlære kviler på bioelektrikk, fri-energi-prinsippet, evolusjonære overgangar i individualitet, like mykje som på Stiny, Hatchuel og Wright. Tverrfagleg er ikkje ei dyd her; det er eit vilkår.

Arkivet er enormt. Det er underutnytta. Det ventar.

\end{foreordbody}
\end{widebody}
